% Options for packages loaded elsewhere
\PassOptionsToPackage{unicode}{hyperref}
\PassOptionsToPackage{hyphens}{url}
\documentclass[
]{article}
\usepackage{xcolor}
\usepackage{amsmath,amssymb}
\setcounter{secnumdepth}{-\maxdimen} % remove section numbering
\usepackage{iftex}
\ifPDFTeX
  \usepackage[T1]{fontenc}
  \usepackage[utf8]{inputenc}
  \usepackage{textcomp} % provide euro and other symbols
\else % if luatex or xetex
  \usepackage{unicode-math} % this also loads fontspec
  \defaultfontfeatures{Scale=MatchLowercase}
  \defaultfontfeatures[\rmfamily]{Ligatures=TeX,Scale=1}
\fi
\usepackage{lmodern}
\ifPDFTeX\else
  % xetex/luatex font selection
\fi
% Use upquote if available, for straight quotes in verbatim environments
\IfFileExists{upquote.sty}{\usepackage{upquote}}{}
\IfFileExists{microtype.sty}{% use microtype if available
  \usepackage[]{microtype}
  \UseMicrotypeSet[protrusion]{basicmath} % disable protrusion for tt fonts
}{}
\makeatletter
\@ifundefined{KOMAClassName}{% if non-KOMA class
  \IfFileExists{parskip.sty}{%
    \usepackage{parskip}
  }{% else
    \setlength{\parindent}{0pt}
    \setlength{\parskip}{6pt plus 2pt minus 1pt}}
}{% if KOMA class
  \KOMAoptions{parskip=half}}
\makeatother
\usepackage{longtable,booktabs,array}
\newcounter{none} % for unnumbered tables
\usepackage{multirow}
\usepackage{calc} % for calculating minipage widths
% Correct order of tables after \paragraph or \subparagraph
\usepackage{etoolbox}
\makeatletter
\patchcmd\longtable{\par}{\if@noskipsec\mbox{}\fi\par}{}{}
\makeatother
% Allow footnotes in longtable head/foot
\IfFileExists{footnotehyper.sty}{\usepackage{footnotehyper}}{\usepackage{footnote}}
\makesavenoteenv{longtable}
\usepackage{graphicx}
\makeatletter
\newsavebox\pandoc@box
\newcommand*\pandocbounded[1]{% scales image to fit in text height/width
  \sbox\pandoc@box{#1}%
  \Gscale@div\@tempa{\textheight}{\dimexpr\ht\pandoc@box+\dp\pandoc@box\relax}%
  \Gscale@div\@tempb{\linewidth}{\wd\pandoc@box}%
  \ifdim\@tempb\p@<\@tempa\p@\let\@tempa\@tempb\fi% select the smaller of both
  \ifdim\@tempa\p@<\p@\scalebox{\@tempa}{\usebox\pandoc@box}%
  \else\usebox{\pandoc@box}%
  \fi%
}
% Set default figure placement to htbp
\def\fps@figure{htbp}
\makeatother
\setlength{\emergencystretch}{3em} % prevent overfull lines
\providecommand{\tightlist}{%
  \setlength{\itemsep}{0pt}\setlength{\parskip}{0pt}}
\usepackage{bookmark}
\IfFileExists{xurl.sty}{\usepackage{xurl}}{} % add URL line breaks if available
\urlstyle{same}
\hypersetup{
  hidelinks,
  pdfcreator={LaTeX via pandoc}}

\author{}
\date{}

\begin{document}

\textbf{\textbar UNIVERSIDAD PERUANA UNIÓN}

\textbf{FACULTAD DE INGENIERÍA Y ARQUITECTURA}


\includegraphics[width=2.3368in,height=1.96176in]{media/image1.png}

\textbf{E.P. INGENIERÍA DE SISTEMAS}

\textbf{Título del proyecto:}

\textbf{``Proyecto de Automatización de una''}

\textbf{Integrantes:}

\textbf{Jhack Ernesto Quispe Mamani}\\
\textbf{Abel Is-Boset Guevara Huasco}\\
\textbf{Javier Kevin Tello Rivas}\\
\textbf{Richard Joseft Yanqui Inocente}\\
\textbf{Alvaro Kenji Cuti Magaño}\\
\textbf{David Guerra Ordoñez}\\
\textbf{Joshua Alfonzo Javier}\\
\textbf{Sebastían Chinchay Collachagua}

\textbf{Lima -- Perú}

\textbf{2025}

\textbf{Índice}

3. Situación de la empresa

3.1. Descripción de la empresa

3.2. Misión

3.3. Visión

3.4. Alineamientos estratégicos (Mapa estrategíco)

3.5. Estructura organizacional

3.6. Enfoque de procesos (Politica de calidad)

3.7. Modelo de Negocio

3. Situación de la empresa

3.1. Descripción de la empresa

La infraestructura tecnológica de Eco Logistic S.A.C. es actualmente
básica y está orientada principalmente a soportar las operaciones
administrativas y de comunicación interna. La organización cuenta con un
acceso a internet centralizado mediante un router, el cual permite la
conectividad de los equipos disponibles dentro de las instalaciones. En
cuanto al equipamiento, dispone de tres computadoras destinadas a
actividades administrativas y operativas, así como una impresora
utilizada para la gestión documental. La comunicación interna y
coordinación operativa se realiza principalmente a través de la
aplicación de mensajería WhatsApp, lo que evidencia la ausencia de
sistemas corporativos integrados de comunicación y gestión. Respecto a
la seguridad, la empresa cuenta con un sistema de videovigilancia
compuesto por 16 cámaras distribuidas en puntos estratégicos de la
infraestructura, las cuales son monitoreadas mediante un televisor que
permite la visualización en tiempo real. Este sistema contribuye al
control de accesos y a la protección de los activos físicos y la
mercadería almacenada. No obstante, la infraestructura de TI presenta
limitaciones significativas en términos de escalabilidad, integración y
automatización de procesos, debido a la falta de sistemas especializados
y plataformas tecnológicas robustas. Esta situación genera una alta
dependencia de procesos manuales y herramientas básicas, lo que impacta
en la eficiencia operativa y en la trazabilidad de la información.

3.2. Misión:

Brindar soluciones logísticas integrales mediante servicios de
almacenamiento, gestión de inventarios y despacho, apoyados en el uso de
tecnología, con el fin de garantizar eficiencia, seguridad y
trazabilidad, generando valor para nuestros clientes.

3.3. Visión:

Ser una empresa logística líder a nivel nacional con proyección
internacional, reconocida por su eficiencia operativa, uso de tecnología
avanzada y capacidad para gestionar terminales de almacenamiento bajo
estándares de calidad y seguridad.

3.4. Alineamientos estratégicos

El plan de acción presentado corresponde a la ejecución anual de los
objetivos estratégicos de Eco Logistic SAC, estructurado bajo un enfoque
de mejora continua y alineado al Cuadro de Mando Integral. Cada objetivo
cuenta con indicadores de seguimiento, responsables designados, un
cronograma mensual y un presupuesto estimado, lo que permite una gestión
organizada, controlada y medible. Entre los objetivos más relevantes se
encuentran la realización de gastos controlados para optimizar costos en
la implementación de soluciones, la inversión en infraestructura de
tecnologías de la información y conectividad para reducir limitaciones
operativas y aumentar la escalabilidad, y la reducción de costos
operativos mediante la adopción de tecnologías de automatización. Estos
objetivos responden directamente a la necesidad de modernización
tecnológica de la empresa, así como a la mejora de la eficiencia en sus
procesos logísticos. El plan establece acciones distribuidas a lo largo
del año, priorizando los meses de mayor impacto operativo, y contempla
un presupuesto total de S/ 77,900. Esta planificación permite alinear
los procesos internos con los resultados esperados, asegurar una
adecuada asignación de recursos tecnológicos y financieros, y facilitar
la toma de decisiones orientadas a fortalecer la competitividad y
sostenibilidad de la organización.

{\def\LTcaptype{none} % do not increment counter
\begin{longtable}[]{@{}
  >{\centering\arraybackslash}p{(\linewidth - 30\tabcolsep) * \real{0.2055}}
  >{\centering\arraybackslash}p{(\linewidth - 30\tabcolsep) * \real{0.1243}}
  >{\centering\arraybackslash}p{(\linewidth - 30\tabcolsep) * \real{0.0881}}
  >{\centering\arraybackslash}p{(\linewidth - 30\tabcolsep) * \real{0.0397}}
  >{\centering\arraybackslash}p{(\linewidth - 30\tabcolsep) * \real{0.0407}}
  >{\centering\arraybackslash}p{(\linewidth - 30\tabcolsep) * \real{0.0409}}
  >{\centering\arraybackslash}p{(\linewidth - 30\tabcolsep) * \real{0.0401}}
  >{\centering\arraybackslash}p{(\linewidth - 30\tabcolsep) * \real{0.0410}}
  >{\centering\arraybackslash}p{(\linewidth - 30\tabcolsep) * \real{0.0405}}
  >{\centering\arraybackslash}p{(\linewidth - 30\tabcolsep) * \real{0.0409}}
  >{\centering\arraybackslash}p{(\linewidth - 30\tabcolsep) * \real{0.0398}}
  >{\centering\arraybackslash}p{(\linewidth - 30\tabcolsep) * \real{0.0399}}
  >{\centering\arraybackslash}p{(\linewidth - 30\tabcolsep) * \real{0.0406}}
  >{\centering\arraybackslash}p{(\linewidth - 30\tabcolsep) * \real{0.0409}}
  >{\centering\arraybackslash}p{(\linewidth - 30\tabcolsep) * \real{0.0409}}
  >{\centering\arraybackslash}p{(\linewidth - 30\tabcolsep) * \real{0.0963}}@{}}
\toprule\noalign{}
\endhead
\bottomrule\noalign{}
\endlastfoot
\multirow{2}{=}{\centering\arraybackslash \textbf{Objetivo Estratégico}}
& \multirow{2}{=}{\centering\arraybackslash \textbf{Indicadores}} &
\multirow{2}{=}{\centering\arraybackslash \textbf{Responsable}} &
\multicolumn{12}{>{\centering\arraybackslash}p{(\linewidth - 30\tabcolsep) * \real{0.4858} + 22\tabcolsep}}{%
\textbf{Cronograma}} &
\multirow{2}{=}{\centering\arraybackslash \textbf{Presupuesto}} \\
& & & \textbf{ene} & \textbf{feb} & \textbf{mar} & \textbf{abr} &
\textbf{may} & \textbf{jun} & \textbf{jul} & \textbf{ago} & \textbf{set}
& \textbf{oct} & \textbf{nov} & \textbf{dic} \\
Realizar gastos controlados para OPTIMIZAR costos en la Implementación
de Soluciones WMS/ERP en la Nube & Rentabilidad sobre Ventas & Jefa de
ventas & X & X & X & & & & & & & & & & \$ 8,500.00 \\
\multirow{2}{=}{\centering\arraybackslash INVERTIR en la infraestructura
TI y conectividad para reducir limitaciones operativas y aumentar la
escalabilidad} & \multirow{2}{=}{\centering\arraybackslash Nivel de
escalabilidad (capacidad instalada vs uso)} &
\multirow{2}{=}{\centering\arraybackslash Gerente general} &
\multirow{2}{=}{\centering\arraybackslash } &
\multirow{2}{=}{\centering\arraybackslash X} &
\multirow{2}{=}{\centering\arraybackslash X} &
\multirow{2}{=}{\centering\arraybackslash X} &
\multirow{2}{=}{\centering\arraybackslash } &
\multirow{2}{=}{\centering\arraybackslash } &
\multirow{2}{=}{\centering\arraybackslash } &
\multirow{2}{=}{\centering\arraybackslash } &
\multirow{2}{=}{\centering\arraybackslash } &
\multirow{2}{=}{\centering\arraybackslash } &
\multirow{2}{=}{\centering\arraybackslash } &
\multirow{2}{=}{\centering\arraybackslash } &
\multirow{2}{=}{\centering\arraybackslash \$ 1,550.00} \\
 \\
REDUCIR costos operativos mediante la adopción de tecnologías de
automatización & Reducción de costos operativos en porcentaje & Gerente
General / Contable & & & & X & X & X & & & & & & & \$ 15,000.00 \\
Ofrecer servicios logísticos personalizados y rápidos usando tecnologías
de automatización & Tiempo de respuesta al cliente; N° de clientes con
trazabilidad en tiempo real & Jefa de ventas & & & & & X & X & X & & & &
& & \$ 5,000.00 \\
Garantizar la seguridad de la información y cumplimiento con las
normativas de protección de datos personales & N° de incidentes de
seguridad; cumplimiento Ley N°29733 (\%) & Analista de inventario & & &
& & & X & X & X & & & & & \$ 3,000.00 \\
Optimizar los procesos logísticos y mejorar la gestión de inventarios
mediante la implementación de WMS/ERP & \% de exactitud de inventario;
N° de transacciones registradas digitalmente & Jefe de operaciones & & &
X & X & X & X & X & & & & & & \$ 10,000.00 \\
Implementar un sistema de gestión de residuos electrónicos y disposición
de equipos en desuso & N° de equipos dados de baja correctamente;
cumplimiento normativo RAEE (\%) & Jefe de Comercialización & & & & & &
& & X & X & & & & \$ 1,000.00 \\
Capacitar al personal en nuevas tecnologías para mejorar la adaptación a
la automatización y la digitalización & N° de empleados capacitados; \%
de reducción de errores operativos post-capacitación & Socia /
Subgerente & & & & & X & X & & & & & & & \$ 6,000.00 \\
Establecer un plan de soporte y mantenimiento preventivo para nuevas
soluciones tecnológicas & N° de mantenimientos realizados; tiempo medio
de resolución de incidencias (horas) & Jefe de supervisión & & & & & & &
& & X & X & X & X & \$ 3,600.00 \\
Fomentar la cultura de innovación y mejora continua a través de
capacitación y gestión del cambio & N° de iniciativas de mejora
propuestas por el equipo; índice de adopción tecnológica (\%) & Socia /
Subgerente & & & & & & & & & & X & X & X & \$ 2,800.00 \\
& & &
\multicolumn{2}{>{\centering\arraybackslash}p{(\linewidth - 30\tabcolsep) * \real{0.0805} + 2\tabcolsep}}{%
} &
\multicolumn{2}{>{\centering\arraybackslash}p{(\linewidth - 30\tabcolsep) * \real{0.0809} + 2\tabcolsep}}{%
} &
\multicolumn{2}{>{\centering\arraybackslash}p{(\linewidth - 30\tabcolsep) * \real{0.0815} + 2\tabcolsep}}{%
} &
\multicolumn{2}{>{\centering\arraybackslash}p{(\linewidth - 30\tabcolsep) * \real{0.0807} + 2\tabcolsep}}{%
} &
\multicolumn{2}{>{\centering\arraybackslash}p{(\linewidth - 30\tabcolsep) * \real{0.0806} + 2\tabcolsep}}{%
} &
\multicolumn{2}{>{\centering\arraybackslash}p{(\linewidth - 30\tabcolsep) * \real{0.0817} + 2\tabcolsep}}{%
} & \textbf{\$ 56,450.00} \\
& & &
\multicolumn{2}{>{\centering\arraybackslash}p{(\linewidth - 30\tabcolsep) * \real{0.0805} + 2\tabcolsep}}{%
} &
\multicolumn{2}{>{\centering\arraybackslash}p{(\linewidth - 30\tabcolsep) * \real{0.0809} + 2\tabcolsep}}{%
} &
\multicolumn{2}{>{\centering\arraybackslash}p{(\linewidth - 30\tabcolsep) * \real{0.0815} + 2\tabcolsep}}{%
} &
\multicolumn{2}{>{\centering\arraybackslash}p{(\linewidth - 30\tabcolsep) * \real{0.0807} + 2\tabcolsep}}{%
} &
\multicolumn{2}{>{\centering\arraybackslash}p{(\linewidth - 30\tabcolsep) * \real{0.0806} + 2\tabcolsep}}{%
} &
\multicolumn{2}{>{\centering\arraybackslash}p{(\linewidth - 30\tabcolsep) * \real{0.0817} + 2\tabcolsep}}{%
} & \textbf{Total} \\
\end{longtable}
}

3.5. Estructura organizacional:

La estructura organizacional de Cologistic S.A.C. se presenta de forma
jerárquica y funcional, distribuida en tres niveles: estratégico,
operativo y de ejecución. Cabe señalar que, debido a que el área de
Tecnologías de la Información cuenta únicamente con un ingeniero
responsable, el presente análisis se centra exclusivamente en el
organigrama del negocio, sin considerar una desagregación estructural
del área de TI. En el nivel estratégico se ubican la Gerencia General,
la Subgerencia y la Gerencia de Contabilidad, responsables de la
dirección global, la administración financiera y la toma de decisiones
clave. El nivel operativo está conformado por la Jefatura de
Operaciones, el Analista de Inventario y la Jefatura de Ventas, quienes
coordinan el control de la trazabilidad y los procesos centrales de la
empresa. Finalmente, en el nivel de ejecución se encuentran los
encargados que realizan las actividades logísticas específicas como la
supervisión de despachos y picking, la operación técnica de estiba, así
como la seguridad armada y control de accesos en garita, permitiendo una
gestión segura, eficiente y alineada con los estándares de sus clientes
corporativos.

\includegraphics[width=5.90625in,height=3.91667in]{media/image3.png}

3.6. Enfoque de procesos

En Cologistic S.A.C., la gestión organizacional se fundamenta en un
enfoque sistémico orientado a la prestación de servicios logísticos
integrales y arrendamiento de espacios a nivel corporativo (B2B).
Entendemos que la confiabilidad de nuestra cadena de suministro no
depende de acciones aisladas, sino de un flujo continuo y altamente
controlado que garantiza la trazabilidad, desde la recepción en nuestros
patios de maniobras hasta el despacho final.

Nuestra arquitectura de procesos se clasifica en tres niveles para
garantizar una operación fluida e ininterrumpida:

\begin{itemize}
\item
  Procesos Estratégicos: Orientan la expansión y sostenibilidad del
  negocio. Incluyen la negociación inmobiliaria para la adquisición de
  nuevas locaciones, la planificación financiera (gestión de contratos
  en dólares) y la gestión comercial B2B.
\item
  Procesos Operativos (Misionales): Constituyen el \emph{core} del
  negocio e incluyen la Recepción y Despacho, la Gestión de Inventarios,
  y los Servicios Logísticos Adicionales (reempaque, repaletizado y
  \emph{picking}).
\item
  Procesos de Soporte: Garantizan la continuidad operativa. Abarcan la
  Gestión de Seguridad (auditorías EPP, personal armado y control de
  plagas) , la Gestión de TI (implementación y mantenimiento de sistemas
  WMS/ERP) y el mantenimiento de infraestructura y maquinaria pesada.
\item
  Política de Calidad y Seguridad Logística Cologistic S.A.C. se
  compromete a brindar soluciones de almacenamiento y gestión logística
  confiables, seguras y adaptables a las necesidades de nuestros
  clientes corporativos y estacionales. Nuestros pilares son:
\end{itemize}

\begin{enumerate}
\def\labelenumi{\arabic{enumi}.}
\item
  Seguridad e Integridad de la Carga: Garantizamos la protección total
  de la mercadería frente a la intemperie y riesgos externos mediante
  infraestructura adecuada y un estricto protocolo de seguridad privada
  ininterrumpida.
\item
  Trazabilidad y Transparencia: Nos comprometemos a mantener una
  co-gestión estrecha con nuestros clientes, asegurando la exactitud de
  los inventarios mediante auditorías físicas recurrentes y reportes
  transparentes de movimientos en tiempo real.
\item
  Prevención y Cumplimiento Normativo: Mantenemos un entorno de trabajo
  seguro auditando rigurosamente el uso de EPP, exigiendo el SCTR para
  todo el personal operativo, y aplicando controles sanitarios estrictos
  para el resguardo de áreas de alimentos
\end{enumerate}

\end{document}
